\section{冬季星空與天體}
\begin{itemize}
    \item 冬季大三角：獵戶座\(\alpha\)（參宿四）、大犬座\(\alpha\)（天狼星） 、小犬座\(\alpha\)（南河三）
    \item 冬季大橢圓：獵戶座\(\beta\)（參宿七）、大犬座\(\alpha\)（天狼星）、小犬座\(\alpha\)（南河三）、雙子座\(\beta\)（北河二）、雙子座\(\alpha\)（北河三）、御夫座\(\alpha\)（五車二）、金牛座\(\alpha\)（畢宿五）
    \item 冬季認星歌訣：
        \begin{center}
            三星高照入寒冬\(\space\)昴星成團亮晶晶\(\space\)金牛低頭衝獵戶\(\space\)群星燦爛放光明 \\
            御夫五星五邊形\(\space\)天河上面放風箏\(\space\)冬夜星空認星座\(\space\)全天最亮天狼星
        \end{center}
    \item 認星方法
        \begin{enumerate}
            \item 冬季大三角：天狼、南河三與參宿四連成一等腰三角形
                \begin{center}
                    獵戶腰帶\(\xrightarrow{\text{東南}}\)天狼星（大犬座）\\
                    雙肩二星\(\xrightarrow{\text{東方}}\)南河三（小犬座）
                \end{center}
            \item 金牛座
                \begin{align*}
                    &\text{天狼星和獵戶的左肩（參宿五）連線}\\
                    &\xrightarrow{\text{西北}}\text{畢宿五（金牛座）（位於牛角Ｖ字型）}\\
                    &\xrightarrow{\text{西北}}\text{昴宿（七姊妹，M45，疏散星團）}
                \end{align*}
            \item 冬季大橢圓
                \begin{align*}
                    &\text{天狼、南河三、參宿七、畢宿五連成一個弧線，再繼續連成一個橢圓}\\
                    &\rightarrow\text{御夫座（五邊形）的五車二}\\
                    &\rightarrow\text{雙子座（像個「北」字型）的北河二與北河三}
                \end{align*}
        \end{enumerate}
\end{itemize}

\subsection{獵戶座 Orion}
\begin{enumerate}
    \item 參宿一、參宿二及參宿三是獵戶腰帶，容易我們判斷獵戶座位子。
    \item 獵戶座方便找尋很多其他星座。（如：大犬座、金牛座、小犬座及雙子座）
    \item 最亮星：參宿七（獵戶座\(\beta\)），全天第7亮星，它的視星等為0.12等。是個三合星的系統，主星為參宿七A藍色超巨星
    \item 參宿四（獵戶座\(\alpha\)）：全天第十亮星，獵戶座第二亮星。紅色半規則變星，視星等在0.2至1.2等之間變化，是一顆紅超巨星，並且是已知第二大和最亮的恆星之一。
\end{enumerate}
獵戶座中的深空天體：

\begin{center}
\begin{longtable}{c | c | c | c | c | c}
    % \hline
    編號 & 種類 & 別名 & 視尺寸 & RA & Dec \\
    \hline
    \endfirsthead

    % \hline
    % 編號 & 種類 & 別名 & 視尺寸 & 拍攝技巧 & 曝光時間 \\
    % \hline
    % \endhead

    % \hline
    % \multicolumn{6}{r}{接下頁} \\
    \endfoot

    % \hline
    \endlastfoot

    M42/NGC 1976 & 瀰漫星雲 & 獵戶座大星雲 & \(\ang{1;18;}\) & \(5^\text{h }35^\text{m }17^\text{s}\) & \(-\ang{5;23;28}\) \\
    M43/NGC 1982 & 發射星雲 & 火鳥星雲 & \(\ang{;20;}\times\ang{;15;}\) & \(5^\text{h }35^\text{m }31^\text{s}\) & \(-\ang{5;16;12}\) \\
    NGC 1980 & 疏散星團 & 獵戶座失落的寶石 & \(\ang{;14;}\) & \(5^\text{h }35^\text{m }24^\text{s}\) & \(-\ang{5;54;54}\) \\
    NGC 1981 & 疏散星團 & - & \(\ang{;28;}\) & \(5^\text{h }35^\text{m }09^\text{s}\) & \(-\ang{4;25;54}\) \\
    NGC 1977 & 反射星雲 & 跑步者星雲 & \(\ang{;10;}\times\ang{;5;}\) & \(5^\text{h }35^\text{m }18^\text{s}\) & \(-\ang{4;41;5}\)\\
    IC 434 & 暗星雲 & 馬頭星雲 & \(\ang{1;30;}\times\ang{;30;}\) & \(5^\text{h }41^\text{m }00^\text{s}\) & \(-\ang{2;30;0}\) \\
    NGC 2024/Sh2-277 & 發射星雲 & 火焰星雲 & \(\ang{;30;}\) & \(5^\text{h }41^\text{m }43^\text{s}\) & \(-\ang{1;50;30}\) \\
    M78/NGC 2068 & \makecell{反射星雲\\（廣域暗星雲）} & \makecell{Casper the Friendly\\Ghost Nebula} & - & \(5^\text{h }46^\text{m }46^\text{s}\) & \(+\ang{0;4;45}\) \\
    NGC 2174 & 發射星雲 & 猴頭星雲 & \(\ang{0;40;}\) & \(6^\text{h }09^\text{m }42^\text{s}\) & \(+\ang{20;30;0}\) \\
    Sh2-261 & 發射星雲 & Lower's Nebula & \(\ang{;30;}\times\ang{;15;}\) & \(6^\text{h }08^\text{m }44^\text{s}\) & \(+\ang{15;40;30}\)
\end{longtable}
\end{center}

\subsection{大犬座 Canis Major}
\begin{enumerate}
    \item 最亮星：天狼星（視星等為−1.46），是夜空中的最亮恆星，放出的光芒約為太陽的20倍
    \item 獵戶的腰帶往東南方延伸能找到天狼星（大犬座\(\alpha\)）
    \item 拉丁語原名意指大狗，代表跟隨天上獵人（獵戶座）的狗。與代表小狗的小犬座對應
    \item 銀河穿過大犬座，星座範圍有許多疏散星團
\end{enumerate}
大犬座中的深空天體：

\begin{center}
\begin{longtable}{c | c | c | c | c | c}
    % \hline
    編號 & 種類 & 別名 & 視尺寸 & RA & Dec \\
    \hline
    \endfirsthead

    % \hline
    % 編號 & 種類 & 別名 & 視尺寸 & 拍攝技巧 & 曝光時間 \\
    % \hline
    % \endhead

    % \hline
    % \multicolumn{6}{r}{接下頁} \\
    \endfoot

    % \hline
    \endlastfoot

    M41/NGC 2287 & 疏散星團 & 小蜂巢星團 & \(\ang{;39;}\) & \(6^\text{h }46^\text{m }01^\text{s}\) & \(-\ang{20;45;24}\) \\
    NGC 2359/Sh2-298 & WR星 & 雷神索爾的頭盔 & \(\ang{;8;}\times\ang{;16;}\) & \(7^\text{h }18^\text{m }30^\text{s}\) & \(-\ang{13;13;48}\) \\
    NGC 2362 & 疏散星團 & 大犬座\(\tau\)星團 & \(\ang{;5;}\) & \(7^\text{h }18^\text{m }41^\text{s}\) & \(-\ang{24;57;18}\) \\
    Sh2-308 & 發射星雲 & 海豚星雲 & \(\ang{;35;}\) & \(6^\text{h }54^\text{m }13^\text{s}\) & \(-\ang{23;55;42}\) \\
\end{longtable}
\end{center}

\subsection{小犬座 Canis Minor}
\begin{enumerate}
    \item 只有兩顆恆星亮度超過4.0視星等
        \begin{itemize}
            \item 最亮星：南河三，0.34視星等，在夜空所有恆星排第八。是一顆白色恆星及一顆白矮星為伴
            \item 南河二，2.9視星等
        \end{itemize}
    \item 拉丁語原名意為小犬，與代表大犬的大犬座對應
    \item 著名的流星雨有：小犬座流星雨、南河三流星雨
\end{enumerate}
小犬座中的深空天體：

\begin{center}
\begin{longtable}{c | c | c | c | c | c}
    % \hline
    編號 & 種類 & 別名 & 視尺寸 & RA & Dec \\
    \hline
    \endfirsthead

    % \hline
    % 編號 & 種類 & 別名 & 視尺寸 & 拍攝技巧 & 曝光時間 \\
    % \hline
    % \endhead

    % \hline
    % \multicolumn{6}{r}{接下頁} \\
    \endfoot

    % \hline
    \endlastfoot

    NGC 2459 & 疏散星團 & - & \(\ang{;2;42}\) & \(7^\text{h }52^\text{m }01^\text{s}\) & \(+\ang{9;33;26}\) \\
    NGC 2508 & 橢圓星系 & - & \(\ang{;2;}\times\ang{;16;}\) & \(8^\text{h }01^\text{m }57^\text{s}\) & \(+\ang{8;33;6}\) \\
    NGC 2402 & 橢圓螺旋星系對 & - & \(\ang{;1;30}\) & \(7^\text{h }30^\text{m }47^\text{s}\) & \(+\ang{9;39;2}\)

\end{longtable}
\end{center}

\subsection{天兔座 Lepus}
\begin{enumerate}
    \item 最亮星：天兔座\(\alpha\)星（廁一），視星等2.6等
    \item 代表著獵人（獵戶座）追逐的兔子
\end{enumerate}
天兔座中的深空天體：

\begin{center}
\begin{longtable}{c | c | c | c | c | c}
    % \hline
    編號 & 種類 & 別名 & 視尺寸 & RA & Dec \\
    \hline
    \endfirsthead

    % \hline
    % 編號 & 種類 & 別名 & 視尺寸 & 拍攝技巧 & 曝光時間 \\
    % \hline
    % \endhead

    % \hline
    % \multicolumn{6}{r}{接下頁} \\
    \endfoot

    % \hline
    \endlastfoot

    M79/NGC 1904 & 球狀星團 & - & \(\ang{;9;36}\) & \(5^\text{h }24^\text{m }11^\text{s}\) & \(-\ang{24;31;27}\) \\
    NGC 2017 & 疏散星團 & - & \(\ang{;5;}\) & \(5^\text{h }39^\text{m }16^\text{s}\) & \(-\ang{17;50;59}\) 
\end{longtable}
\end{center}

\subsection{麒麟座 Monoceros}
\begin{enumerate}
    \item 最亮星：麒麟座\(\alpha\)星，3.9等
\end{enumerate}
麒麟座中的深空天體：

\begin{center}
\begin{longtable}{c | c | c | c | c | c}
    % \hline
    編號 & 種類 & 別名 & 視尺寸 & RA & Dec \\
    \hline
    \endfirsthead

    % \hline
    % 編號 & 種類 & 別名 & 視尺寸 & 拍攝技巧 & 曝光時間 \\
    % \hline
    % \endhead

    % \hline
    % \multicolumn{6}{r}{接下頁} \\
    \endfoot

    % \hline
    \endlastfoot

    M50/NGC 2323 & 疏散星團 & - & \(\ang{;15;}\) & \(7^\text{h }02^\text{m }48^\text{s}\) & \(-\ang{8;20;16}\) \\
    NGC 2232 & 疏散星團 & - & \(\ang{;10;}\) & \(6^\text{h }28^\text{m }00^\text{s}\) & \(-\ang{4;50;51}\) \\
    NGC 2264 & \makecell{星團\\（廣域發射星雲）} & 聖誕樹星雲 & \makecell{\(\ang{;20;}\)\\（星雲\(\ang{1;30;}\)）} & \(6^\text{h }40^\text{m }58^\text{s}\) & \(+\ang{9;53;43}\) \\
    NGC 2237 & 發射星雲 & 玫瑰星雲 & \(\ang{1;20;}\times\ang{;50;}\) & \(6^\text{h }30^\text{m }55^\text{s}\) & \(+\ang{5;2;57}\) \\
    IC 2177 & 反射/發射星雲 & 海鷗星雲 & \(\ang{2;;}\times\ang{;40;}\) & \(7^\text{h }04^\text{m }25^\text{s}\) & \(-\ang{10;27;18}\)
\end{longtable}
\end{center}

\subsection{金牛座 Taurus}
\begin{enumerate}
    \item 金牛座就在獵戶座的右上方，獵戶的腰帶向西北方延伸就會碰到畢宿五（金牛座\(\alpha\)）
    \item 最亮星：畢宿五（金牛座\(\alpha\)），視星等不規則地在0.75 等到 0.95 等之變化，全天第十四亮星，紅巨星
    \item 北半球冬季夜空上最大、最顯著的星座之一
    \item 畢宿亮星排列稱V字形結構，又稱為金牛座V字
    \item 五車五（金牛座\(\beta\)）：第二亮星，全天第27亮星，視星等1.65
    \item 著名的流星雨有：金牛座南流星雨、金牛座北流星雨
\end{enumerate}
金牛座中的深空天體：

\begin{center}
\begin{longtable}{c | c | c | c | c | c}
    % \hline
    編號 & 種類 & 別名 & 視尺寸 & RA & Dec \\
    \hline
    \endfirsthead

    % \hline
    % 編號 & 種類 & 別名 & 視尺寸 & 拍攝技巧 & 曝光時間 \\
    % \hline
    % \endhead

    % \hline
    % \multicolumn{6}{r}{接下頁} \\
    \endfoot

    % \hline
    \endlastfoot

    M45 & 疏散星團 & 昴宿星團 & \(\ang{1;30;}\) & \(3^\text{h }46^\text{m }38^\text{s}\) & \(+\ang{24;10;41}\) \\
    M1/NGC 1952 & 超新星殘骸 & 蟹狀星雲 & \(\ang{;6;}\times\ang{;4;}\) & \(5^\text{h }34^\text{m }32^\text{s}\) & \(+\ang{22;0;52}\) \\
    NGC 1555 & \makecell{反射（變光）星雲\\廣域暗星雲} & 欣德變光星雲 & \(\ang{;2;}\) & \(4^\text{h }21^\text{m }59^\text{s}\) & \(+\ang{19;32;06}\) \\
    NGC 1514 & 行星狀星雲 & 水晶球星雲 & \(\ang{;2;12}\) & \(4^\text{h }09^\text{m }16^\text{s}\) & \(+\ang{30;46;33}\)
\end{longtable}
\end{center}

\subsection{御夫座 Auriga}
\begin{enumerate}
    \item 最亮星：御夫座\(\alpha\)（五車二），分光雙星，0.1 等。是全天第六亮星，肉眼觀測略呈黃色，含有兩顆黃色巨星
    \item 風箏狀五角形的御夫座可用來尋找北極星（面向北方，沿著五角形逆時針方向數，在五車二的東邊或左手邊分別為五車三和五車四。將五車四向五車三方向延伸，約5倍遠的距離即為北極星）
    \item 拉丁語原名意為「戰車御者」，即駕駛戰車的人
    \item 御夫座流星雨：又稱「御夫座\(\alpha\)流星雨」，以間歇爆發聞名。每年約8月28日開始，9月1日清晨達到巔峰
\end{enumerate}
御夫座中的深空天體：

\begin{center}
\begin{longtable}{c | c | c | c | c | c}
    % \hline
    編號 & 種類 & 別名 & 視尺寸 & RA & Dec \\
    \hline
    \endfirsthead

    % \hline
    % 編號 & 種類 & 別名 & 視尺寸 & 拍攝技巧 & 曝光時間 \\
    % \hline
    % \endhead

    % \hline
    % \multicolumn{6}{r}{接下頁} \\
    \endfoot

    % \hline
    \endlastfoot

    M36/NGC 1960 & 疏散星團 & Pinwheel Cluster & \(\ang{;10;}\) & \(5^\text{h }36^\text{m }18^\text{s}\) & \(+\ang{34;8;24}\) \\
    M37/NGC 2099 & 疏散星團 & \makecell{Salt and Pepper\\Cluster} & \(\ang{;24;}\) & \(5^\text{h }52^\text{m }18^\text{s}\) & \(+\ang{32;33;12}\) \\
    M38/NGC 1912 & 疏散星團 &  海星星團 & \(\ang{;20;}\) & \(5^\text{h }28^\text{m }43^\text{s}\) & \(+\ang{35;51;18}\) \\
    NGC 1907 & 疏散星團 & - & \(\ang{;7;}\) & \(5^\text{h }28^\text{m }05^\text{s}\) & \(+\ang{35;19;30}\) \\
    IC 410/NGC 1893 & 發射星雲 & 蝌蚪星雲 & \(\ang{;40;}\times\ang{;30;}\) & \(5^\text{h }22^\text{m }44^\text{s}\) & \(+\ang{33;24;43}\) \\
    NGC 2281 & 疏散星團 & \makecell{Broken Heart\\Cluster} & \(\ang{;25;}\) & \(6^\text{h }48^\text{m }17^\text{s}\) & \(+\ang{41;4;42}\) \\
    NGC 1931 & 發射星雲 & 蒼蠅星雲 & \(\ang{;7;}\) & \(5^\text{h }31^\text{m }23^\text{s}\) & \(+\ang{34;13;59}\) \\
    NGC 1664 & 疏散星團 & - & \(\ang{;11;24}\) & \(4^\text{h }51^\text{m }05^\text{s}\) & \(+\ang{43;40;34}\) \\
    NGC 1778 & 疏散星團 & - & \(\ang{;4;30}\) & \(5^\text{h }08^\text{m }05^\text{s}\) & \(+\ang{37;1;22}\) \\
    NGC 1857 & 疏散星團 & - & \(\ang{;4;30}\) & \(5^\text{h }20^\text{m }05^\text{s}\) & \(+\ang{39;20;36}\) \\
    IC 405 & 發射/反射星雲 & 火焰之星星雲 & \(\ang{;50;}\times\ang{;30;}\) & \(5^\text{h }16^\text{m }28^\text{s}\) & \(+\ang{34;21;22}\) \\
    IC 417 & 發射星雲 & 蜘蛛星雲 & \(\ang{;13;}\times\ang{;10;}\) & \(5^\text{h }28^\text{m }05^\text{s}\) & \(+\ang{34;25;26}\)
\end{longtable}
\end{center}

\subsection{雙子座 Gemini}
\begin{enumerate}
    \item 最亮星：北河三（雙子座\(\beta\)），視星等為1.14，橘色巨星，是全天第十七亮星
    \item 其中兩顆亮星雙子座\(\alpha\)跟\(\beta\)（北河二和北河三）為雙子座兩兄弟的頭部
    \item 在雙子座H附近發現天王星，在雙子座\(\delta\)附近發現冥王星
    \item 北河二（雙子座\(\alpha\)），是第一顆經過天文觀測確認的物理雙星，複合星等1.58，全天第22亮星。子星\(\alpha\)A的視星等是1.9等，子星\(\alpha\)B是2.9等，兩顆藍白色恆星形成物理雙星，環繞週期約 470 年。此外還有一個視星等為8.8等的伴星\(\alpha\)C，而且這三顆星分別都是分光雙星。因此北河二是一顆六合星
    \item 雙子座流星雨:是三大流星雨之一每年的高峰期大約都是在12月13、14日左右
\end{enumerate}
雙子座中的深空天體：

\begin{center}
\begin{longtable}{c | c | c | c | c | c}
    % \hline
    編號 & 種類 & 別名 & 視尺寸 & RA & Dec \\
    \hline
    \endfirsthead

    % \hline
    % 編號 & 種類 & 別名 & 視尺寸 & 拍攝技巧 & 曝光時間 \\
    % \hline
    % \endhead

    % \hline
    % \multicolumn{6}{r}{接下頁} \\
    \endfoot

    % \hline
    \endlastfoot

    M35/NGC 2168 & 疏散星團 & - & \(\ang{;40;}\) & \(6^\text{h }08^\text{m }54^\text{s}\) & \(+\ang{24;20;0}\) \\
    NGC 2392 & 行星狀星雲 & 愛斯基摩星雲 & \(\ang{;;20}\) & \(7^\text{h }29^\text{m }11^\text{s}\) & \(+\ang{20;54;42}\) \\
    Sh2-274 & 行星狀星雲 & 蛇妖星雲 & \(\ang{;10;15}\) & \(7^\text{h }29^\text{m }03^\text{s}\) & \(+\ang{13;14;48}\) \\
    IC 444 & 反射星雲 & - & \(\ang{;8;}\times\ang{;4;}\) & \(6^\text{h }18^\text{m }33^\text{s}\) & \(+\ang{23;18;47}\) \\
    IC 443 & 超新星殘骸 & 水母星雲 & \(\ang{;50;}\times\ang{;40;}\) & \(6^\text{h }16^\text{m }36^\text{s}\) & \(+\ang{22;31;54}\)
\end{longtable}
\end{center}

\subsection{巨蟹座Cancer}
\begin{enumerate}
    \item 是最黯淡的黃道星座，可以透過其他星座的位子確定其所在（東邊為獅子頭部，西北是雙子的北河二和北河三，西南是小犬的南河三）
    \item 最亮星：柳宿增十（巨蟹座\(\beta\)），視星等為3.52
\end{enumerate}
巨蟹座中的深空天體：

\begin{center}
\begin{longtable}{c | c | c | c | c | c}
    % \hline
    編號 & 種類 & 別名 & 視尺寸 & RA & Dec \\
    \hline
    \endfirsthead

    % \hline
    % 編號 & 種類 & 別名 & 視尺寸 & 拍攝技巧 & 曝光時間 \\
    % \hline
    % \endhead

    % \hline
    % \multicolumn{6}{r}{接下頁} \\
    \endfoot

    % \hline
    \endlastfoot

    M44/NGC 2632 & 疏散星團 & 蜂巢星團 & \(\ang{1;10;}\) & \(8^\text{h }40^\text{m }24^\text{s}\) & \(+\ang{19;40;0}\) \\
    M67/NGC 2682 & 疏散星團 & - & \(\ang{;25;}\) & \(8^\text{h }51^\text{m }18^\text{s}\) & \(+\ang{11;48;0}\) \\
    NGC 2775 & 螺旋星系 & - & \(\ang{;4;15}\) & \(9^\text{h }10^\text{m }19^\text{s}\) & \(+\ang{7;2;16}\) \\
    \makecell{NGC 2535\\NGC2536} & \makecell{螺旋星系\\棒旋星系} & （正在交互作用） & \(\ang{;2;30}\) & \(8^\text{h }11^\text{m }12^\text{s}\) & \(+\ang{25;12;24}\) \\
\end{longtable}
\end{center}

\subsection{鹿豹座 Camelopardalis}
\begin{enumerate}
    \item 是一個黯淡的星座，組成鹿豹座的星星都是暗星
    \item 在日本稱為麒麟座。「鹿豹」指的是長頸鹿
\end{enumerate}
鹿豹座中的深空天體：

\begin{center}
\begin{longtable}{c | c | c | c | c | c}
    % \hline
    編號 & 種類 & 別名 & 視尺寸 & RA & Dec \\
    \hline
    \endfirsthead

    % \hline
    % 編號 & 種類 & 別名 & 視尺寸 & 拍攝技巧 & 曝光時間 \\
    % \hline
    % \endhead

    % \hline
    % \multicolumn{6}{r}{接下頁} \\
    \endfoot

    % \hline
    \endlastfoot

    NGC 1502 & 疏散星團 & - & \(\ang{;10;12}\) & \(4^\text{h }07^\text{m }49^\text{s}\) & \(+\ang{62;19;53}\) \\
    NGC 2403 & 中間棒旋星系 & - & \(\ang{;21;54}\times\ang{;12;18}\) & \(7^\text{h }36^\text{m }51^\text{s}\) & \(+\ang{65;36;9}\) \\
    NGC 342 & 中間棒旋星系 & - & \(\ang{;20;}\) & \(3^\text{h }46^\text{m }49^\text{s}\) & \(+\ang{68;5;46}\)
\end{longtable}
\end{center}

\subsection{天貓座 Lynx}
\begin{enumerate}
    \item 最亮星：天貓座\(\alpha\)星（軒轅四），3.1等
    \item 拉丁語原名意指山貓，意指未來只有長著山貓般利眼的後輩才能看到
\end{enumerate}
天貓座中的深空天體：

\begin{center}
\begin{longtable}{c | c | c | c | c | c}
    % \hline
    編號 & 種類 & 別名 & 視尺寸 & RA & Dec \\
    \hline
    \endfirsthead

    % \hline
    % 編號 & 種類 & 別名 & 視尺寸 & 拍攝技巧 & 曝光時間 \\
    % \hline
    % \endhead

    % \hline
    % \multicolumn{6}{r}{接下頁} \\
    \endfoot

    % \hline
    \endlastfoot

    NGC 2419 & 球狀星團 & 星系浪子 & \(\ang{;4;30}\) & \(7^\text{h }38^\text{m }08^\text{s}\) & \(+\ang{38;52;47}\) \\
    NGC 2537 & 棒旋星系 & 熊掌星系 & \(\ang{;2;}\) & \(8^\text{h }13^\text{m }14^\text{s}\) & \(\ang{45;59;23}\) \\
    \makecell{NGC 2841\\NGC 2541\\NGC 2500\\NGC 2552} & 星系群 & NGC 2841星系群 & - & \makecell{Ursa Major\\\(8^\text{h }14^\text{m }40^\text{s}\)\\\(8^\text{h }01^\text{m }53^\text{s}\)\\\(8^\text{h }19^\text{m }20^\text{s}\)} & \makecell{Ursa Major\\\(+\ang{49;03;41}\)\\\(+\ang{50;44;14}\)\\\(+\ang{50;0;21}\)} \\
    NGC 2683 & 螺旋星系 & - & \(\ang{;10;}\times\ang{;3;}\) & \(8^\text{h }52^\text{m }41^\text{s}\) & \(+\ang{33;25;18}\) \\
    JnEr 1 & 行星狀星雲 & 耳機星雲 & \(\ang{;6;40}\) & \(7^\text{h }57^\text{m }52^\text{s}\) & \(+\ang{53;25;17}\) \\
\end{longtable}
\end{center}

\subsection{劍魚座 Dorado}
\begin{enumerate}
    \item 南半球星座
    \item 最亮星：劍魚座\(\alpha\)星，視星等3.3等
    \item 象徵海洋生物劍魚
    \item 大麥哲倫雲主要位於此星座內，是距離銀河系最近的星系
    \item 超新星1987A最亮達2.8等1604年以來地球所見最亮的超新星。在1987年爆發，位於大麥哲倫雲內接近蜘蛛星雲的地方。肉眼可見的期間長達 10 個月之久。
\end{enumerate}
劍魚座中的深空天體：

\begin{center}
\begin{longtable}{c | c | c | c | c | c}
    % \hline
    編號 & 種類 & 別名 & 視尺寸 & RA & Dec \\
    \hline
    \endfirsthead

    % \hline
    % 編號 & 種類 & 別名 & 視尺寸 & 拍攝技巧 & 曝光時間 \\
    % \hline
    % \endhead

    % \hline
    % \multicolumn{6}{r}{接下頁} \\
    \endfoot

    % \hline
    \endlastfoot

    ESO056-115 & 棒旋星系 & 大麥哲倫星系（LMC） & \(\ang{10;15}\times\ang{9;10}\) & \(5^\text{h }23^\text{m }35^\text{s}\) & \(-\ang{69;45;21}\)\\
    NGC 1763 & 發射星雲 & （屬於LMC） & \(\ang{;5;12}\times\ang{;3;18}\) & \(4^\text{h }56^\text{m }48^\text{s}\) & \(-\ang{64;24;33}\) \\
    NGC 2070 & 發射星雲 & 蜘蛛星雲（LMC） & \(\ang{;16;}\) & \(5^\text{h }38^\text{m }41^\text{s}\) & \(-\ang{69;6;2}\)\\
    NGC 1566 & 中間螺旋星系 & - & \(\ang{;7;}\times\ang{;5;}\) & \(7^\text{h }38^\text{m }08^\text{s}\) & \(+\ang{38;52;47}\) \\
\end{longtable}
\end{center}

\subsection{船帆座 Vela}
\begin{enumerate}
    \item 由天舟座一分為三
\end{enumerate}
船帆座中的深空天體：
\begin{center}
\begin{longtable}{c | c | c | c | c | c}
    % \hline
    編號 & 種類 & 別名 & 視尺寸 & RA & Dec \\
    \hline
    \endfirsthead

    % \hline
    % 編號 & 種類 & 別名 & 視尺寸 & 拍攝技巧 & 曝光時間 \\
    % \hline
    % \endhead

    % \hline
    % \multicolumn{6}{r}{接下頁} \\
    \endfoot

    % \hline
    \endlastfoot

    NGC 2547 & 疏散星團 & - & \(\ang{;7;48}\) & \(8^\text{h }10^\text{m }09^\text{s}\) & \(-\ang{49;12;20}\) \\
    NGC 3132 & 行星狀星雲 & 八裂星雲 & \(\ang{;;30}\) & \(10^\text{h }07^\text{m }01^\text{s}\) & \(-\ang{40;26;12}\) \\
    IC 2391 & 疏散星團 & \makecell{Omicron Velae\\Cluster} & \(\ang{;29;}\) & \(8^\text{h }40^\text{m }32^\text{s}\) & \(-\ang{53;2;8}\)
\end{longtable}
\end{center}

\subsection{船底座 Carina}
\begin{enumerate}
    \item 最亮星：船底座\(\alpha\)星（老人星）白色超巨星，-0.6 等全天第二亮恆星
\end{enumerate}
船底座中的深空天體：

\begin{center}
\begin{longtable}{c | c | c | c | c | c}
    % \hline
    編號 & 種類 & 別名 & 視尺寸 & RA & Dec \\
    \hline
    \endfirsthead

    % \hline
    % 編號 & 種類 & 別名 & 視尺寸 & 拍攝技巧 & 曝光時間 \\
    % \hline
    % \endhead

    % \hline
    % \multicolumn{6}{r}{接下頁} \\
    \endfoot

    % \hline
    \endlastfoot

    NGC 2516 & 疏散星團 & - & \(\ang{;24;18}\) & \(7^\text{h }58^\text{m }06^\text{s}\) & \(-\ang{60;45;12}\) \\
    NGC 3114 & 疏散星團 & - & \(\ang{;12;18}\) & \(10^\text{h }02^\text{m }30^\text{s}\) & \(-\ang{60;7;50}\) \\
    NGC 3372 & 發射星雲 & 船底座星雲 & \(\ang{2;;}\) & \(10^\text{h }45^\text{m }08^\text{s}\) & \(-\ang{59;52;0}\) \\
    NGC 3532 & 疏散星團 & 許願井星團 & \(\ang{;12;}\) & \(11^\text{h }05^\text{m }48^\text{s}\) & \(-\ang{58;46;14}\) \\
    IC 2602 & 疏散星團 & 南昴宿星團 & \(\ang{;48;}\) & \(10^\text{h }42^\text{m }56^\text{s}\) & \(-\ang{64;23;39}\) \\
\end{longtable}
\end{center}

\subsection{天鴿座 Columba}
\begin{enumerate}
    \item 最亮星：天鴿座\(\alpha\)星 (丈人一)，藍白色恆星，視星等2.6 等
\end{enumerate}
天鴿座中的深空天體：

\begin{center}
\begin{longtable}{c | c | c | c | c | c}
    % \hline
    編號 & 種類 & 別名 & 視尺寸 & RA & Dec \\
    \hline
    \endfirsthead

    % \hline
    % 編號 & 種類 & 別名 & 視尺寸 & 拍攝技巧 & 曝光時間 \\
    % \hline
    % \endhead

    % \hline
    % \multicolumn{6}{r}{接下頁} \\
    \endfoot

    % \hline
    \endlastfoot

    NGC 1851 & 球狀星團 & - & \(\ang{;9;}\) & \(5^\text{h }14^\text{m }07^\text{s}\) & \(-\ang{40;02;48}\) \\
\end{longtable}
\end{center}

\subsection{船尾座 Puppis}
\begin{enumerate}
    \item 最亮星：船尾座\(\zeta\)星（弧矢二十二）2.2 等
    \item 由天舟座一分為三，船尾座是其中最大的星座
\end{enumerate}
船尾座中的深空天體：

\begin{center}
\begin{longtable}{c | c | c | c | c | c}
    % \hline
    編號 & 種類 & 別名 & 視尺寸 & RA & Dec \\
    \hline
    \endfirsthead

    % \hline
    % 編號 & 種類 & 別名 & 視尺寸 & 拍攝技巧 & 曝光時間 \\
    % \hline
    % \endhead

    % \hline
    % \multicolumn{6}{r}{接下頁} \\
    \endfoot

    % \hline
    \endlastfoot

    M46/NGC 2437 & 疏散星團 & - & \(\ang{;21;}\) & \(7^\text{h }41^\text{m }47^\text{s}\) & \(-\ang{14;48;36}\) \\
    M93/NGC 2447 & 疏散星團 & - & \(\ang{;15;}\) & \(7^\text{h }44^\text{m }29^\text{s}\) & \(-\ang{23;51;11}\) \\
    NGC 2451 & 疏散星團 & - &  \(\ang{;37;30}\) & \(7^\text{h }45^\text{m }15^\text{s}\) & \(-\ang{37;58;3}\) \\
    NGC 2477 & 疏散星團 & - & \(\ang{;18;48}\) & \(7^\text{h }52^\text{m }09^\text{s}\) & \(-\ang{38;31;59}\) \\
    M47/NGC 2422 & 疏散星團 & - & \(\ang{;19;48}\) & \(7^\text{h }36^\text{m }34^\text{s}\) & \(-\ang{14;28;56}\) \\
\end{longtable}
\end{center}

\subsection{繪架座 Pictor}
\begin{enumerate}
    \item 最亮星：繪架座\(\alpha\)，是一顆白色主序星
    \item 繪架座是一個黯淡的星座；它最亮的三顆恆星都位於顯著的老人星的附近
\end{enumerate}
繪架座中的深空天體：
\begin{center}
\begin{longtable}{c | c | c | c | c | c}
    % \hline
    編號 & 種類 & 別名 & 視尺寸 & RA & Dec \\
    \hline
    \endfirsthead

    % \hline
    % 編號 & 種類 & 別名 & 視尺寸 & 拍攝技巧 & 曝光時間 \\
    % \hline
    % \endhead

    % \hline
    % \multicolumn{6}{r}{接下頁} \\
    \endfoot

    % \hline
    \endlastfoot

    NGC 1705 & 不規則矮星系 & - & \(\ang{;2;}\) & \(4^\text{h }54^\text{m }14^\text{s}\) & \(-\ang{53;21;39}\) \\
\end{longtable}
\end{center}

\subsection{雕具座 Caelum}
\begin{enumerate}
    \item 最亮星：雕具座\(\alpha\)星，4.4 等
\end{enumerate}

\subsection{時鐘座 Horologium}
\begin{enumerate}
    \item 南天球六顆黯淡恆星組成的星座最亮的天園增六亮度達到3.85，是老化的橙巨星
    \item 時鐘座R是米拉變星，視星等在4.7至14.3範圍變化
\end{enumerate}
時鐘座中的深空天體：

\begin{center}
\begin{longtable}{c | c | c | c | c | c}
    % \hline
    編號 & 種類 & 別名 & 視尺寸 & RA & Dec \\
    \hline
    \endfirsthead

    % \hline
    % 編號 & 種類 & 別名 & 視尺寸 & 拍攝技巧 & 曝光時間 \\
    % \hline
    % \endhead

    % \hline
    % \multicolumn{6}{r}{接下頁} \\
    \endfoot

    % \hline
    \endlastfoot

    NGC 1261 & 球狀星團 & - & \(\ang{;5;6}\) & \(3^\text{h }12^\text{m }15^\text{s}\) & \(-\ang{55;13;0}\) \\
    NGC 1512 & 棒旋星系 & - & \(\ang{;8;25}\times\ang{;4;}\) & \(4^\text{h }03^\text{m }54^\text{s}\) & \(-\ang{43;20;56}\) 
\end{longtable}
\end{center}

\subsection{波江座 Eridanus}
\begin{enumerate}
    \item 面積排名第六
    \item 最亮星：波江座\(\alpha\)星（水委一）藍白色恆星，0.5 等，全天第九亮星，位於波江座最南端。
\end{enumerate}
波江座中的深空天體：

\begin{center}
\begin{longtable}{c | c | c | c | c | c}
    % \hline
    編號 & 種類 & 別名 & 視尺寸 & RA & Dec \\
    \hline
    \endfirsthead

    % \hline
    % 編號 & 種類 & 別名 & 視尺寸 & 拍攝技巧 & 曝光時間 \\
    % \hline
    % \endhead

    % \hline
    % \multicolumn{6}{r}{接下頁} \\
    \endfoot

    % \hline
    \endlastfoot

    NGC 1909/IC 2118 & 反射星雲 & 女巫頭星雲 & \(\ang{3;;}\times\ang{1;;}\) & \(5^\text{h }04^\text{m }54^\text{s}\) & \(-\ang{7;15;56}\) \\
\end{longtable}
\end{center}

\subsection{網罟座 Reticulum}
\begin{enumerate}
    \item 最亮星：網罟座\(\alpha\)星，3.3 等
\end{enumerate}